\begin{abstract}
\noindent
Do paragliders move differently when they fly in a group? Group time differs from solo flight in
many ways, but so do the pilots who fly in groups, the places and hours where groups form, and the
moments pilots choose to fly together. We measure the motion by the local exponent $\Hloc(\tau)$ of
the time-averaged mean squared displacement and by statistics of the climbs and glides, and take
each of the other causes away with one test: contrasts within strata of site, day and hour and
within the same flight; solo flights reweighted to the zones and hours of the groups; the same route
flown alone on another day; the same flight before, during and after a group. Over a quarter of an hour a pilot in a group progresses
straighter. $\Delta\Hloc(1000~\mathrm{s})$ is \DHlocThousandFlight{} against the same flight alone and
\DHlocThousandGroupSdh{} against the same flights alone at the same site, day and hour, while who
flies in groups accounts for \DHlocThousandSel. Reweighting the solo flights to the zones and hours
where groups fly lowers the long-lag gap by at most a fifth; on the half of the group time that can
be matched, fixing the route lowers it by about \RouteExplained\%; and on the same flight the
straighter progress starts and stops with the group. Within one glide or climb the progress is no straighter. Glides are faster and flatter
in a group, both within strata and on the same flight before and after; pilots also climb more, but on the same flight the climbs change when the
group forms and not when it breaks. Most of this is there with two pilots, and the pilot ahead ends
searches in a climb as often as the one behind. The tests fix the pilot, the place, the route, the
hour and the moment, but not the conditions of the day that bring pilots together: the rest of the
effect goes with flying together, and whether the group causes it stays open.
\end{abstract}
